\chapter*{Conclusions}

\addcontentsline{toc}{chapter}{Conclusions}

We proved that Polynomial Calculus with Resolution is analyzable, and hence so are Nullstellensatz and Polynomial Calculus, by lifting the block-width argument to block-degree. As a corollary, we obtained a new proof that automating PCR is $\mathsf{NP}$-hard. We then compared analyzability with feasible interpolation. Proving that interpolation implies analyzability would show that Resolution is not weakly automatable unless $\mathsf{P}$ = $\mathsf{NP}$, while for every disjoint $\mathsf{NP}$-pair there is an analyzable system, simulating Resolution and closed under restrictions, that has feasible interpolation only if the pair is polynomial-time separable. Any connection between the two notions must therefore use specific properties of natural proof systems. Several questions that further research may explore remain open.

\textit{Search-analyzability of PCR}. The most immediate question left open is whether a satisfying assignment of $\varphi$ can be extracted in polynomial time from a short PCR proof of $\lnot\mathrm{Ref}_s(\varphi)$.

\textit{Pudlák's upper bound in algebraic systems}. Our proof that automating PCR is $\mathsf{NP}$-hard uses Pudlák's upper bound for Resolution, which PCR inherits because it simulates Resolution. Nullstellensatz and Polynomial Calculus do not simulate Resolution efficiently, so the same argument does not give the $\mathsf{NP}$-hardness of automating them, even though both are analyzable by our result. It is open whether $\mathrm{Ref}_s(\varphi)$, or a suitable algebraic variant of it, has short Polynomial Calculus or Nullstellensatz refutations when $\varphi$ is satisfiable. A positive answer, combined with our analyzability results, would give a direct proof that automating these systems is $\mathsf{NP}$-hard along the lines of Atserias and Müller.

\textit{Analyzability of Cutting Planes, Res(k) and bounded-depth Frege}. Every analyzability result needs a lower bound for $\mathrm{Ref}_s(\varphi)$ in the analyzed system. For Res(k), the block-width technique may extend through switching lemmas for k-DNFs. For Cutting Planes, the usual techniques are lifting from communication complexity and monotone interpolation, and it is unclear whether either adapts to $\mathrm{Ref}_s(\varphi)$. For bounded-depth Frege, switching-lemma lower bounds for the pigeonhole principle exist, but the retraction version would need a separate analysis.

\textit{A natural separation of analyzability and interpolation}. We showed that analyzability and feasible interpolation are different in general, but the separating system is artificial. Is there a natural analyzable proof system without feasible interpolation, or a natural system with feasible interpolation that is not analyzable? Cutting Planes is the natural candidate for the second question, since it has feasible interpolation through monotone real circuits and its analyzability is open. A proof that Cutting Planes is not analyzable under a standard assumption would give the first natural separation. A proof that it is analyzable would extend the known coincidence of the two properties to a system that is incomparable with PCR.

\textit{Does interpolation imply analyzability on restricted systems?} Proposition 5.6 rules out a general proof that feasible interpolation implies analyzability, unless the weak automatability of Resolution is settled. It does not rule out such a proof for restricted classes of systems. A natural candidate is the class of systems that extend Resolution, are closed under restrictions and in which lines are evaluated by a bounded number of Boolean functions, such as Res(k). A result of this kind would explain why every analyzable system known before this thesis also has feasible interpolation.

\textit{Weak automatability of Resolution}. The relationship between analyzability and interpolation is tied to the weak automatability of Resolution, one of the central open problems of the area. Weak automatability asks for an algorithm that finds a proof in some possibly stronger system, so the meta-complexity method does not apply directly: the stronger system may refute $\mathrm{Ref}_s(\varphi)$ efficiently even when $\varphi$ is unsatisfiable. Any progress on analyzability for systems above Resolution, such as Res(2), therefore bears directly on this problem.